\section{Цель работы}

Ознакомиться с основными методами формирования заданных интервалов времени и перевода микроконтроллера в режим пониженного энергопотребления.
Закрепить навыки работы с осциллографом и оценочной платой MCBSTM32F200 в качестве измерительного генератора.

\section{Постановка задачи}

Разработать программу для микроконтроллера STM32F200, которая при помощи таймера формирует периодическое попеременное включение и выключение светодиодов PG6 и PG7 с заданными временными характеристиками.
При нажатии на кнопку WAKEUP программа должна переводить микроконтроллер в спящий режим, а при ее отпускании --- пробуждать микроконтроллер.

\section{Выполнение работы}

В среде Keil $\mu$Vision5 был создан проект, в котором реализовано управление светодиодами с использованием таймера и обработка внешнего прерывания.
Ниже приведены ключевые фрагменты программы и пояснения.

\subsection{Функция задержки}

Для создания программной задержки используется простая пустая петля.
Переменная цикла объявлена с квалификатором \verb|volatile|, чтобы компилятор не оптимизировал цикл.

\begin{listing}[H]
\caption{Delay function (milliseconds)}
\begin{minted}[linenos]{c}
void delay_ms(uint32_t milliseconds) {
    volatile uint32_t i;
    for (i = 0; i < (milliseconds * 4200U); i++) { __NOP(); }
}
\end{minted}
\end{listing}

\subsection{Инициализация GPIO}

Включается тактирование портов GPIOA и GPIOG, настраиваются выводы PG6 и PG7 как выходы, а PA0 как вход для кнопки WAKEUP.

\begin{listing}[H]
\caption{GPIO initialization}
\begin{minted}[linenos]{c}
RCC->AHB1ENR |= RCC_AHB1ENR_GPIOAEN | RCC_AHB1ENR_GPIOGEN;

/* PG6, PG7 */
GPIOG->MODER &= ~((3U << (6*2)) | (3U << (7*2)));
GPIOG->MODER |=   (1U << (6*2)) | (1U << (7*2));

/* PA0 (WAKEUP) */
GPIOA->MODER &= ~(3U << (0*2));
\end{minted}
\end{listing}

\subsection{Настройка таймера TIM2}

Таймер TIM2 настраивается с предделителем и автоперезагрузкой.
После настройки включается прерывание по обновлению, задается приоритет обработчика и разрешается работа таймера.

\begin{listing}[H]
\caption{TIM2 setup}
\begin{minted}[linenos]{c}
RCC->APB1ENR |= RCC_APB1ENR_TIM2EN;

TIM2->PSC = 119990U;
TIM2->ARR = 249U;

TIM2->CR1  |= TIM_CR1_ARPE | TIM_CR1_URS;
TIM2->DIER |= TIM_DIER_UIE;
TIM2->CR1  |= TIM_CR1_CEN;

NVIC_SetPriority(TIM2_IRQn, 0);
NVIC_EnableIRQ(TIM2_IRQn);
\end{minted}
\end{listing}

\subsection{Обработчик прерывания таймера}

При каждом переполнении таймера выполняется переключение светодиодов PG6 и PG7.
После обработки события сбрасывается флаг прерывания таймера.

\begin{listing}[H]
\caption{TIM2 IRQ Handler}
\begin{minted}[linenos]{c}
void TIM2_IRQHandler(void) {
    if (TIM2->SR & TIM_SR_UIF) {
        GPIOG->ODR ^= (1U << 6) | (1U << 7);
        TIM2->SR &= ~TIM_SR_UIF;
    }
}
\end{minted}
\end{listing}

\subsection{Настройка внешнего прерывания EXTI0 (WAKEUP)}

Включается тактирование SYSCFG, линия EXTI0 привязывается к PA0, разрешается прерывание по обоим фронтам и настраивается NVIC.

\begin{listing}[H]
\caption{EXTI0 setup}
\begin{minted}[linenos]{c}
RCC->APB2ENR |= RCC_APB2ENR_SYSCFGEN;

SYSCFG->EXTICR[0] &= ~SYSCFG_EXTICR1_EXTI0;
SYSCFG->EXTICR[0] |=  SYSCFG_EXTICR1_EXTI0_PA;

EXTI->IMR  |= EXTI_IMR_MR0;
EXTI->RTSR |= EXTI_RTSR_TR0;
EXTI->FTSR |= EXTI_FTSR_TR0;

NVIC_SetPriority(EXTI0_IRQn, 0);
NVIC_EnableIRQ(EXTI0_IRQn);
\end{minted}
\end{listing}

\subsection{Обработчик EXTI0}

Обработчик реализует подавление дребезга, переключение режимов сна (\verb|SLEEPDEEP|, \verb|SLEEPONEXIT|) и изменение типа срабатывания EXTI для корректного входа в сон и пробуждения.

\begin{listing}[H]
\caption{EXTI0 IRQ Handler}
\begin{minted}[linenos]{c}
void EXTI0_IRQHandler(void) {
    if (EXTI->PR & EXTI_PR_PR0) {
        delay_ms(100); /* debounce */

        if (!sleep_mode) {
            /* Enter deep sleep on exit from IRQ */
            SCB->SCR |= SCB_SCR_SLEEPDEEP_Msk;
            SCB->SCR |= SCB_SCR_SLEEPONEXIT_Msk;

            /* Switch EXTI to falling edge for wake-up */
            EXTI->RTSR &= ~EXTI_RTSR_TR0;
            EXTI->FTSR |= EXTI_FTSR_TR0;
            sleep_mode = true;
        } else {
            /* Exit sleep mode */
            SCB->SCR &= ~SCB_SCR_SLEEPONEXIT_Msk;
            SCB->SCR &= ~SCB_SCR_SLEEPDEEP_Msk;

            /* Restore EXTI to trigger on rising edge */
            EXTI->RTSR |= EXTI_RTSR_TR0;
            EXTI->FTSR &= ~EXTI_FTSR_TR0;
            sleep_mode = false;
        }

        /* Clear EXTI pending flag */
        EXTI->PR = EXTI_PR_PR0;
    }
}
\end{minted}
\end{listing}

\section{Работа с осциллографом}

При исследовании сигнала на выходе микроконтроллера был получен прямоугольный сигнал.
По экрану осциллографа период сигнала составил примерно $T = 1.9$ с, что соответствует частоте около $f = 0.53$ Гц.
По встроенным измерениям осциллографа были определены времена нарастания и спада сигнала, после чего фронт и спад были рассмотрены на растянутой временной развертке.

\screenshot{IMG_20260414_123652.jpg}{Осциллограмма прямоугольного сигнала на выходе микроконтроллера.}

На фотографии \ref{fig:IMG_20260414_123907.jpg} показан результат автоматического измерения времени нарастания.
Измеренное время нарастания составило $t_r = 10.16~\text{нс}$, размах сигнала составил примерно $V_{pp}=3.64$ В.

\screenshot{IMG_20260414_123907.jpg}{Явное измерение времени нарастания сигнала.}

На фотографии \ref{fig:IMG_20260414_125842.jpg} показан результат автоматического измерения времени спада.
Измеренное время спада равно $t_f = 10.11~\text{нс}$, размах сигнала составил примерно $V_{pp}=3.60$ В.

\screenshot{IMG_20260414_125842.jpg}{Явное измерение времени спада сигнала.}

После явного измерения был растянут фронт сигнала для визуальной проверки формы переходного процесса.

\screenshot{IMG_20260414_125223_1.jpg}{Растянутое изображение фронта сигнала.}

Также был растянут спад сигнала.
На экране осциллографа дополнительно отображалось значение $t_f = 10.10~\text{нс}$.

\screenshot{IMG_20260414_125909.jpg}{Растянутое изображение спада сигнала.}

\section{Вывод}

В ходе выполнения лабораторной работы была разработана программа для микроконтроллера STM32F200, реализующая формирование временных интервалов при помощи таймера TIM2 и обработку внешнего прерывания от кнопки WAKEUP.
Были изучены настройка таймера, обработка его прерывания, работа с EXTI и перевод микроконтроллера в спящий режим.
С помощью осциллографа были измерены параметры выходного сигнала: период около $1.9$ с, частота около $0.53$ Гц, размах около $3.64$ В, время нарастания $10.16$ нс и время спада $10.11$ нс.
